\section{The original carrier from the near field to radiation}
\label{sec:actual-radiation-zone}
The user-supplied review derived the exact decomposition and compact-edge
remainder below. An independent calculation from the retained complete kernel
agrees. The flux estimate and finite-endpoint examples extend that derivation.
The original six signed $g_{ab}$, equal pair weights, both tensor polarizations,
$\kappa=4/5$, radius balls and $\Delta=2/573049$ remain fixed.

Put $a=R/c_*$, $v=r/R\geq1$, $t_r=t-r/c_*$ and
$x_r=R\phi(r,t)$. Polynomial expansion gives
\begin{align}
 P_2\!\left[-w+\frac{1-w^2}{2v}\right]
 &=P_2(w)-\frac{3w(1-w^2)}{2v}
             +\frac{3(1-w^2)^2}{8v^2},\\
 x_r&=x_{\rm rad}+x_{\rm ind}+x_{\rm near},\nonumber\\
 x_{\rm rad}&=\frac1{\kappa v}\int_{-1}^1P_2(w)f(t_r-aw)\,dw,\nonumber\\
 x_{\rm ind}&=-\frac3{2\kappa v^2}\int_{-1}^1w(1-w^2)f(t_r-aw)\,dw,\nonumber\\
 x_{\rm near}&=\frac3{8\kappa v^3}\int_{-1}^1(1-w^2)^2f(t_r-aw)\,dw.
 \label{eq:exact-three-radial-pieces}
\end{align}
These are exact $r^{-1}$, $r^{-2}$ and $r^{-3}$ pieces at fixed retarded time.
All have the original causal support $[(r-R)/c_*,(r+R)/c_*]$.
The unchanged homogeneous incoming field is added separately. The normalized
angular tensor remains $\sqrt{5/(8\pi)}\Lambda_n$ applied to the original
plus/cross source; its sphere isometry fixes the angular amplitude ruler.

\subsection{A whole-pulse derivative and remainder theorem}
For $f\in C^3(\mathrm R)$ with bounded fourth weak derivative, parity and
Taylor's integral remainder yield
\begin{align}
 x_r(t)&=\frac{2R^3}{15\kappa}
 \left[\frac{f''(t_r)}{c_*^2r}
       +\frac{3f'(t_r)}{c_*r^2}+\frac{3f(t_r)}{r^3}\right]+E(t),\\
 |E(t)|&\leq\frac{a^4\|f^{(4)}\|_\infty}{60\kappa v}
       +\frac{a^3\|f'''\|_\infty}{35\kappa v^2}
       +\frac{a^2\|f''\|_\infty}{35\kappa v^3}.
 \label{eq:compact-edge-radiation-remainder}
\end{align}
For example, the radiative coefficient uses
$\int P_2=0$ and $\int w^2P_2=4/15$; its remainder follows from
$\int |P_2|w^4\leq2/5$. The induction and near remainders use respectively
$\int w^4(1-w^2)=4/35$ and
$\int w^2(1-w^2)^2=16/105$. These estimates apply on the zero extension,
including intervals that cross either pulse edge.

The actual maintained pulse is $f=g_{ab}F(\Delta t)$, with
$F(\theta)=\sin^4(\theta/8)e^{-i\theta}$ on $[0,8\pi]$ and zero elsewhere.
Its first three derivatives match zero at both endpoints. Its five-frequency
expansion gives
\[
 \begin{gathered}
 \|f^{(j)}\|_\infty\leq |g_{ab}|\Delta^jD_j,\\
 (D_0,D_1,D_2,D_3,D_4)
   =\left(1,1,\frac{17}{16},\frac{19}{16},\frac{709}{512}\right).
 \end{gathered}
\]
Define $z=\Delta R/c_*$, $u=\Delta r/c_*$ and the specified positive ruler
$A_0=2R^3|g_{ab}|\Delta^2/(15\kappa c_*^2r)$. Then, for every time,
\begin{equation}
 \begin{aligned}
 &\frac{|x_r-A_0\operatorname{sgn}(g_{ab})F''(\Delta t_r)|}{A_0}\\
 &\quad\leq\frac3u+\frac3{u^2}\\
 &\qquad+z^2\left(\frac{709}{4096}+\frac{57}{224u}
                               +\frac{51}{224u^2}\right).
 \end{aligned}
 \label{eq:radiation-ruler-bound}
\end{equation}
This is an absolute curve bound; phase and relative error need a nonzero
comparator. The full original radius ball gives $z<5.708\times10^{-7}$.
For $u=100,1000,10^6$, the right side is below
$0.030300000001$, $0.003003000001$ and $0.0000030000031$.
These correspond to $r=u/\Delta$ in the same $c_*=1$ model chart.
The old receiver $r=R+64$ has $u<0.000224$.

At the retarded envelope peak $\theta=4\pi$, $F''=-17/16$ is real.
The leading radiative peak therefore has a $6.25\%$ amplitude correction,
with the same phase, relative to a central-frequency unchanged-envelope
approximation. Increasing $r$ leaves that waveform correction present.
The four-cycle maintained pulse and the one-cycle autonomous clock retain
their different source protocols.

\subsection{A finite-radius amplitude ruler and positive emitted energy}
The exact monochromatic identity is
\begin{align}
 \chi_r&=-\frac{2j_2(z)e^{iu}}{\kappa v}
                 \left(1+\frac{3i}{u}-\frac3{u^2}\right),\\
 \chi_{\rm rad}&=-\frac{2j_2(z)e^{iu}}{\kappa v},\\
 \left|\frac{\chi_r}{\chi_{\rm rad}}\right|^2
 &=1+\frac3{u^2}+\frac9{u^4},\\
 \left|\frac{\chi_r}{\chi_{\rm rad}}-1\right|
 &\leq\frac3{|u|}+\frac3{u^2}.
\end{align}
The ratio requires $j_2(z)\ne0$. At $u=100$ the amplitude excess is below
$0.000150045$; the induction term still contributes an order $1/u$ phase.
The local detector functional, angular projection, gain, original source
vertex and $1/r$ factor all enter the recorded amplitude ruler.

The exact coefficient at null infinity and its energy are
\begin{align}
 \psi(t_r)&=\lim_{r\to\infty}r\phi(r,t_r+r/c_*)
       =\frac1\kappa\int_{-1}^1P_2(w)f(t_r-aw)\,dw,\\
 W&=\frac{\kappa}{2c_*}\int_{\mathrm R}|\psi'|^2dt_r.
\end{align}
For $\psi_0=2a^2f''/(15\kappa)$, the comparator is
\[
 W_0=\frac{2R^4g_{ab}^2\Delta^5I_3}{225\kappa c_*^5}.
\]
% Insert within the existing original-carrier section, after W and W_0 are
% defined in radiation_local.tex, replacing the conservative energy paragraph.
\subsection{Exact smoothing fixes the integrated energy sign}
\label{subsec:signed-radiative-energy-update}
The owner's six-page review supplies a sharper theorem for the same
maintained pulse. It keeps both original compact edges and all six signed
source vertices. Put $a=R/c_*$ and let
\[
 \mu(w)=\begin{cases}
  \frac{15}{16}(1-w^2)^2,&|w|\leq1,\\
  0,&|w|>1.
 \end{cases}
\]
Its mass, mean and variance are $1$, $0$ and $1/7$. Since
$P_2(w)=[(1-w^2)^2]''/8$, the two endpoint terms vanish under two
integrations by parts; the equivalent spherical Bessel normalization is also given by DLMF 10.54.1.~[20] Thus the \emph{exact} outgoing radiation is
\[
 \psi=\frac{2a^2}{15\kappa}\mathcal C_\mu f'',\qquad
 \psi'=\frac{2a^2}{15\kappa}\mathcal C_\mu f''',
\]
where $(\mathcal C_\mu g)(t)=\int\mu(w)g(t-aw)\,dw$.
For any $g$ with $g,g'\in L^2$, polarization gives
\begin{align*}
 \|g\|_2^2-\|\mathcal C_\mu g\|_2^2
 &=\frac12\iint\mu(w)\mu(v)\\[-2pt]
 &\quad\times\|g(\cdot-aw)-g(\cdot-av)\|_2^2\,dw\,dv\\
 &\leq\frac{a^2}{7}\|g'\|_2^2.
\end{align*}
The last inequality uses the Sobolev translation bound and
$\int\!\int\mu(w)\mu(v)(w-v)^2\,dw\,dv=2/7$.
The original zero-extended pulse is $C^3$ across its endpoints, so
$g=f'''$ has weak derivative $g'=f^{(4)}\in L^2$. Its original
five-frequency expansion gives
\[
 \begin{aligned}
 I_3&=\int_0^{8\pi}|F'''|^2d\theta=\frac{14245\pi}{4096},\\
 I_4&=\int_0^{8\pi}|F^{(4)}|^2d\theta=\frac{324817\pi}{65536}.
 \end{aligned}
\]
Consequently the same $W,W_0$ satisfy
\begin{align}
 0\leq1-\frac W{W_0}
 &\leq\frac{a^2\|f^{(4)}\|_2^2}{7\|f'''\|_2^2}
   =\frac{324817}{1595440}z^2\nonumber\\
 &<6.633\times10^{-14},\qquad z=\Delta R/c_*.
 \label{eq:signed-energy-deficit-update}
\end{align}
The numerical upper endpoint uses the full original radius ball.
Finite-source smoothing therefore reduces this integrated energy relative
to the derivative comparator. It is a signed certificate for the original
maintained pulse; it changes neither the autonomous-clock preparation nor
its source-depletion account. The common ratio retains each lane's signed
waveform and cancels only its common squared vertex in $W/W_0$.

\subsection{Exact source work, sphere flux and near-field storage}
The supplied review asks for an all-time conservation check for the maintained
incremental pulse. A direct calculation from the original kernel gives it.
Let $a=R/c_*$ and $\mu(w)=15(1-w^2)^2/16$ on $[-1,1]$, and set
\[
 A(u)=\frac{2R^2}{15\kappa c_*^2}
       \int\mu(w)f(u-aw)\,dw,\qquad u=t-r/c_*.
\]
Two integrations by parts give the exact exterior solution
\[
 \phi(r,t)=\frac{A''(u)}r+\frac{3c_*A'(u)}{r^2}
                       +\frac{3c_*^2A(u)}{r^3}.
\]
This includes both compact-pulse edges. Write $q_j=A^{(j)}(u)$ and
$y=(q_2,c_*q_1/r,c_*^2q_0/r^2)$. The complete outward sphere flux is
\[
 \begin{gathered}
 F_r(t)=\frac{\kappa}{2c_*}|q_3|^2+\frac{dS_r(u)}{dt},\\
 S_r=\frac{\kappa}{2r}y^*G y,\qquad
 G=\begin{pmatrix}7/2&6&9/2\\6&15&27/2\\9/2&27/2&27/2\end{pmatrix}.
 \end{gathered}
\]
Its exact LDL pivots are $7/2,33/7,27/44$, all positive. Thus $S_r$ is
nonnegative stored near-field energy on the outgoing characteristic; its
derivative can have either sign. The instantaneous sphere flux need not be
positive. The full exterior energy is
\[
 E_{>r}(t)=\frac{\kappa}{2c_*}\int_{-\infty}^{u}|A'''(v)|^2dv+S_r(u).
\]
Since $A$ is compactly supported, every exterior sphere has the same
all-time integrated flux $W=\kappa\int|A'''|^2/(2c_*)$.

The original radial energy density and flux are
\[
 \begin{gathered}
 e=\frac\kappa4\left[\frac{r^2|\phi_t|^2}{c_*^2}
                     +r^2|\phi_r|^2+6|\phi|^2\right],\\
 F_r=-\frac\kappa2r^2\operatorname{Re}(\phi_t^*\phi_r).
 \end{gathered}
\]
Their local continuity equation and the original source jump give
$\dot E_{<b}+F_b=\operatorname{Re}(f^*\phi_t(R))$ for $b>R$.
Keep $H_{\rm drive}=-\operatorname{Re}(f^*\phi(R))$ and
$H_{\rm ct}=|f|^2/(4R)$ in the actual tensor sector. Explicit control power is
\[
 P_{\rm ctrl}=-\operatorname{Re}(\dot f^*\phi(R))+\dot H_{\rm ct}.
\]
Hence
\[
 \frac d{dt}(E_{<b}+H_{\rm drive}+H_{\rm ct})+F_b=P_{\rm ctrl}.
\]
The maintained source's even energy is constant; any different source-energy
change must be added to its recipient and power account. With zero incoming
field and zero force endpoints, the completed net control work equals $W$,
the integrated flux through every $b>R$, and the full outgoing exterior
energy. Intermediate interaction and counterterm changes remain explicit.
Signed endpoint cancellation supplies no zero-apparatus-cost assertion.

For nonzero incoming field, retain the incoming, emitted and bilinear cross
terms in both energy and flux, including
$\operatorname{Re}(f^*\phi_{{\rm in},t}(R))$ in source power. The frozen native storage check retains the cross block for two outgoing
jets. A separate exact bidirectional checker uses the inward exterior
solution $\phi_{\rm in}=A_{\rm in}''/r-3c_*A_{\rm in}'/r^2+
3c_*^2A_{\rm in}/r^3$ and verifies pure and cross continuity and signed
radiative flux. A regular incoming background also has its matching
outgoing component; an inward characteristic alone is only local to the exterior. This maintained classical carrier
identity remains distinct from autonomous quantum-source depletion.

The exact flux, full outgoing exterior, incoming preparation, source work,
boundary interaction, counterterm and even source/reference energy remain in
the ledger. This fixed flat-carrier normalization supplies no cosmological
$Q$, redshift-dependent generation function or SI conversion by itself.

\section{The canonical electromagnetic comparison}
\label{sec:maxwell-canonical-comparator}
For one transverse conformal Maxwell mode use
$L_M=(A'^2-k^2A^2)/2$. The earlier $Q=1$ comparator uses $h=A/b$ and
$L_h=b^2(h'^2-k^2h^2)/2$. With $\mathcal H=b'/b$,
\begin{align}
 L_h&=\frac12\left[A'^2-
              \left(k^2-\frac{b''}{b}\right)A^2\right]
                   -\frac12(\mathcal H A^2)',\\
 J_hh&=\frac{J_h}{b}A,\qquad j_A=\frac{j_h}{b_E},\\
 D_hh_R&=\frac{D_h}{b_R}A_R.
\end{align}
For the displayed linear $b$, $b''=0$, so these are the same free equations
with the stated canonical boundary term, source and detector transformations.
Their energies satisfy
\[
 \begin{gathered}
 E_h-E_M=-\mathcal H AA'+\tfrac12\mathcal H^2A^2,\\
 E_h=\tfrac12[(A'-\mathcal H A)^2+k^2A^2].
 \end{gathered}
\]
An impulse from zero field has the common work
$j_h^2/(2b_E^2)=j_A^2/2$, recovering the earlier impulse numbers under this
protocol. At nonzero incoming $A$, the two impulse work changes differ by
$-\mathcal H A j_A$; that term belongs to the boundary/background account.
A fixed Maxwell impulse or detector gain requires the corresponding
transformed $h$ values. This specifies the electromagnetic comparison of
the linear-$b$ example.

\section{Moving endpoints of the actual travel-time map}
\label{sec:moving-travel-time-gradient}
Let a differentiable one-branch travel-time function define actual events by
\[
 F(t_E)-t_E=T[t_E,x_E(t_E),F(t_E),x_R(F(t_E))].
\]
Implicit differentiation yields the review's event Jacobian, including
explicit reception-time dependence when present:
\begin{equation}
 F'=\frac{1+\partial_ET+\nabla_ET\cdot v_E}
          {1-\partial_RT-\nabla_RT\cdot v_R}.
 \label{eq:full-moving-event-jacobian}
\end{equation}
The chosen forward causal branch has a nonzero denominator and $F'>0$.
For stationary isotropic optics, endpoint variation gives
$\nabla_ET=-p_E$, $\nabla_RT=p_R$, with optical slowness
$p=n_{\rm opt}\widehat t/c_*$.
The velocities and gradients use their different actual event epochs.
With worldline proper rates $a_A=d\tau_A/dt$ and intrinsic interacting-clock
rates $\nu_A(\tau_A)$, the observed frequency comparison is
\[
 \mathcal R=\frac{a_E\nu_E(\tau_E)}{a_RF'\nu_R(\tau_R)}.
\]
The interacting transient and proper ages stay in this expression. The
actual delay derivative along events is $dT/dt_E=F'-1$.

For the retained weak-source, finite-endpoint logarithm let
$D=|x_R-x_E|$, $A=r_E+r_R$, $n=(x_R-x_E)/D$,
$e_A=x_A/r_A$ and $K=(1+\gamma)GM/c_*^3$. Then
\begin{align}
 T_{\rm app}&=D/c_*+K\log\frac{A+D}{A-D},\\
 \nabla_ET_{\rm app}
 &=-n/c_*-\frac{2K(De_E+An)}{A^2-D^2},\\
 \nabla_RT_{\rm app}
 &= n/c_*+\frac{2K(An-De_R)}{A^2-D^2}.
\end{align}
These gradients enter the actual implicit map. The travel-time approximation
retains its metric-order error. More quantitatively, write its numerator and
denominator as $A_*,B_*$. If $B_{\rm app}\geq b>\delta_B$,
$|A_*-A_{\rm app}|\leq\delta_A$ and
$|B_*-B_{\rm app}|\leq\delta_B$, then
\[
 |F'-F'_{\rm app}|\leq
 \frac{\delta_A}{b-\delta_B}
 +\frac{|A_{\rm app}|\delta_B}{b(b-\delta_B)}.
\]
The operands include path-gradient and endpoint-velocity uncertainties.
Compose outgoing, proper-wait and return maps at their actual event epochs.

\section{A signed finite-frequency transport correction}
For the retained Volterra equation $v=v_0+Vv$ put $v_1=Vv_0$ and
$M=\int|U|/k$. Ordered integration gives
$\|V^nv_0\|_\infty\leq M^n/n!$, hence
\[
 v=v_0+v_1+R_2,\qquad |R_2|\leq e^M-1-M.
\]
The same remainder controls the derivative divided by $k$. A nonzero lower
bound for $|v_0+v_1|$ turns this disk into amplitude and phase bounds about
the signed first correction. In the retained $M=1/5500$ example,
$|R_2|\leq M^2/[2(1-M)]<1.654\times10^{-8}$.
The source-derived physical transport coefficient remains its own calculation.
